\section{Sección 4: Gestión de Proyectos, MRP y Modelamiento (MILP)}

\begin{tcolorbox}[colback=blue!5!white,colframe=blue!60!black,title=Parámetros del Modelo]
\textbf{Variables y Notación:}\\
PERT: $t_o$ (optimista), $t_m$ (probable), $t_p$ (pesimista).\\
MRP (Silver-Meal): $K$ (Costo de emisión/setup), $h$ (Costo de mantener inventario).\\
MILP: $X_{ij}$ (Variable continua de decisión), $Y_i$ (Variable binaria de activación).
\end{tcolorbox}

\subsection{Algoritmo 1: Ruta Crítica y Acortamiento (PERT/CPM)}
Procedimiento para la evaluación de proyectos:

\begin{enumerate}
    \item \textbf{Red PERT y Holguras:} 
    \begin{itemize}
        \item Calcula $ES$ (inicio temprano), $EF=ES+TE$, $LF$ (término tardío) y $LS=LF-TE$.
        \item $\text{Holgura}=LS-ES=LF-EF$. Si la holgura es cero, la actividad pertenece a la
        \textbf{ruta crítica}.
    \end{itemize}
    \item \textbf{Cálculo de Varianzas:} 
    $$ \text{Varianza de una Actividad} = \sigma^2 = \left( \frac{Pesimista - Optimista}{6} \right)^2 $$
    \item \textbf{Probabilidad de Terminar (Z):}
    \begin{itemize}
        \item Varianza de la Ruta ($\sigma^2_{ruta}$) = \textbf{SUMA} de las varianzas de las actividades críticas. \textbf{OJO:} Se suman varianzas, ¡NUNCA se suman las desviaciones estándar!
        \item Para un plazo contractual $T_c$:
        $Z = \frac{T_c - \mu_{ruta}}{\sqrt{\sigma^2_{ruta}}}$ y
        $P(T\le T_c)=\Phi(Z)$.
    \end{itemize}
    \item \textbf{Crashing (Acortamiento):} 
    $$ \text{Costo Marginal} = \frac{Costo_{Acelerado} - Costo_{Normal}}{Tiempo_{Normal} - Tiempo_{Acelerado}} $$
    Siempre acelera la actividad de la ruta crítica con menor costo marginal,
    pero solo mientras tenga margen de reducción. Si aparecen dos rutas críticas,
    debes reducir ambas simultáneamente (o escoger una actividad común) para
    acortar el proyecto.
\end{enumerate}

\vspace{0.4cm}
\begin{ejercicio}{Crashing PERT/CPM (Examen 2023 - Parte 1, Pregunta 2)}
\textbf{Enunciado:} Te dan actividades A, B, C, D, E, F con tiempos normales y de crashing.
\textbf{Resolución Crashing (Paso a Paso):} 1) Identificas la ruta crítica (ej. A-C-E). 2) Divides la diferencia de costo normal/acelerado por los días ahorrados para cada una. 3) Si A cuesta \$100/día y C cuesta \$50/día, \textbf{aceleras C} primero. ¡Nunca aceleres actividades fuera de la ruta crítica porque gastas dinero y el proyecto no se acorta!
\end{ejercicio}

\begin{formulas}{Lectura rápida de una red PERT}
En una pasada hacia adelante, una actividad con varios predecesores comienza
cuando termina el último: $ES_j=\max_i EF_i$. En la pasada hacia atrás, su
término tardío queda determinado por el sucesor más exigente:
$LF_i=\min_j LS_j$. No sumes desviaciones estándar: bajo independencia,
$\sigma_P^2=\sum_{i\in\text{ruta crítica}}\sigma_i^2$ y
$\sigma_P=\sqrt{\sigma_P^2}$. Si el plazo está bajo la media, $Z<0$ y la
probabilidad de cumplimiento es menor que $50\%$.
\end{formulas}

\subsection{Algoritmo 2: MRP y Loteo (Silver-Meal)}
Tienes demandas semanales y un costo de pedir (setup $K$) y mantener ($h$).
Primero explota el BOM y calcula los requerimientos brutos del componente a
partir de los \emph{lanzamientos} planificados de sus padres; no uses
directamente las recepciones del producto padre.

\begin{enumerate}
    \item \textbf{Regla del Costo Promedio (Silver-Meal):} Se van agregando demandas al lote de hoy, hasta que el Costo Promedio por Período empiece a AUMENTAR.
    \item \textbf{Paso 1 (Lote para 1 semana):} Costo Total = $K$. Costo Promedio = $K / 1$.
    \item \textbf{Paso 2 (Lote para 2 semanas):} Costo Total = $K + (Demanda_{sem2} \cdot h \cdot 1)$. Costo Prom = Total / 2.
    \item \textbf{Paso 3 (Lote para 3 semanas):} Costo Total = Anterior $+ (Demanda_{sem3} \cdot h \cdot 2)$. Costo Prom = Total / 3.
    \item \textbf{Corte:} Si el Costo Prom del Paso 3 $>$ Costo Prom del Paso 2, detente. El pedido será solo para las semanas 1 y 2. Inicia un nuevo pedido en la semana 3.
    \item \textbf{MRP completo:} La recepción planificada se ubica en el período de necesidad y
    el lanzamiento se desplaza hacia atrás según el lead time:
    $PORelease_{t-L}=POR_t$. Antes de aplicar Silver-Meal descuenta inventario
    disponible y recepciones programadas.
\end{enumerate}

\vspace{0.4cm}
\begin{ejercicio}{Algoritmo Silver-Meal (Examen 2022 / 2024)}
\textbf{Enunciado Loteo:} Setup $K = \$50$. Costo inventario $h = \$2$/semana. Demanda: Sem1=10, Sem2=15, Sem3=5.
\textbf{Resolución Silver-Meal:}
- \textbf{Opción 1 (Pedir solo S1):} Costo Total = 50. Costo Prom = 50 / 1 = \textbf{50}.
- \textbf{Opción 2 (Pedir S1 y S2):} Pido 25 hoy. Guardo 15 para la otra semana. Costo Total = $50 + (15 \times 2 \times 1) = 80$. Costo Prom = 80 / 2 = \textbf{40}.
- \textbf{Opción 3 (Pedir S1, S2, S3):} Pido 30. Costo = $50 + 30 + (5 \times 2 \times 2) = 100$. Costo Prom = 100 / 3 = \textbf{33.3}.
Sigo agregando mientras el Costo Promedio siga bajando.
\end{ejercicio}

\begin{formulas}{Wagner--Whitin: óptimo global}
Cuando el enunciado pida el óptimo exacto, Silver--Meal es solo una heurística.
Con $f(t)$ igual al costo mínimo acumulado hasta el período $t$:
\[
f(t)=\min_{1\le j\le t}\left\{f(j-1)+K+
H\sum_{r=j}^{t}(r-j)D_r\right\},\qquad f(0)=0.
\]
La alternativa $j$ representa emitir en $j$ una orden que cubre las demandas
$j,\ldots,t$. La política óptima se recupera siguiendo hacia atrás las
alternativas que alcanzan el mínimo.
\end{formulas}

\subsection{Algoritmo 3: Modelamiento MILP (Diagnóstico y Formulación)}
Si te piden escribir la función objetivo (F.O.) y las restricciones.
\begin{enumerate}
    \item \textbf{Variables de Decisión:} $X_{ij}$ (Cantidad producida/enviada), $Y_i$ (Binaria 1 si abro planta, 0 si no).
    \item \textbf{Función Objetivo (Minimizar Costos):} Suma del (Costo Unitario $\cdot X_{ij}$) + (Costo Fijo $\cdot Y_i$).
    \item \textbf{Restricción de Demanda:} Todo lo que entra al cliente debe ser $\ge$ Demanda.
    \item \textbf{Restricción Lógica (Capacidad) - ¡La más importante!:}
    $$ \sum X_{ij} \le M \cdot Y_i $$
    (Si $Y_i=0$, obliga a que $X=0$. Si $Y_i=1$, permite operar hasta $M$).
\end{enumerate}

\begin{formulas}{Plantilla completa de localización capacitada}
Para plantas $i\in I$ y clientes $j\in J$, define $x_{ij}\ge0$ como el flujo
enviado y $y_i\in\{0,1\}$ como la apertura de la planta. Con demanda
$D_j$, capacidad $U_i$, costo fijo $F_i$ y costo unitario $c_{ij}$:
\[
\begin{aligned}
\min\quad &\sum_{i\in I}F_i y_i+\sum_{i\in I}\sum_{j\in J}c_{ij}x_{ij}\\
\text{s.a.}\quad &\sum_{i\in I}x_{ij}\ge D_j &&\forall j\in J,\\
&\sum_{j\in J}x_{ij}\le U_i y_i &&\forall i\in I,\\
&x_{ij}\ge0,\quad y_i\in\{0,1\}.
\end{aligned}
\]
Si hay inventario por período, agrega el balance
$I_{t}=I_{t-1}+P_t-D_t$; si hay activación consecutiva, usa una binaria
$z_t\ge y_t-y_{t-1}$ y cobra el costo fijo sobre $z_t$, no sobre todos los
períodos operados.
\end{formulas}

\vspace{0.4cm}
\begin{ejercicio}{Modelamiento MILP (Examen 2021 / 2023 - Parte 2)}
\textbf{Enunciado MILP:} Tienes 3 plantas posibles para abastecer 5 ciudades. Abrir una planta cuesta $F_i$. El transporte cuesta $C_{ij}$. Plantea el modelo de minimización.
\textbf{Resolución (Variables Clave):}
Escribe esto en tu cuaderno:
Binaria $Y_i = 1$ si abro la planta $i$ y continua $X_{ij}$ igual a la
cantidad enviada de la planta $i$ a la ciudad $j$. La función objetivo es
\[
\min Z=\sum_iF_iY_i+\sum_i\sum_j C_{ij}X_{ij},
\]
acompañada por demanda, capacidad y dominio de las variables. Sin estas
restricciones, la expresión anterior no es todavía un modelo MILP completo.
\end{ejercicio}
